\section{Definite Aliasing Analysis}\label{sec:aliasing}

A \emph{path} is either an expression rooted at a local variable, that follows
\<final> field dereferences or list accesses, or an integer constant.
Paths have a type and a length. A constant $\kappa\in\nat$ limits the maximal
length of a path. We assume that $\kappa$ is a global constant in this work,
so that we not need to clutter each definition with this parameter.
This limit is essential to make our analysis finitely computable, as shown later.
Larger $\kappa$ lead to a more precise analysis,
in principle, but also potentially more expensive.
Therefore, $\kappa$ acts as a lever that selects the desired precision for the
analysis. For simplicity, in the examples that follow, we assume $\kappa=4$.

A relevant path has type list. While paths will be used
to define the definite aliasing analysis, relevant paths will be used in the
subsequent boundness analysis.
%
\begin{definition}[Paths]\label{def:path}
  Given a method $C.m$, the finite set $\paths{C.m}$ of the
  \emph{paths in $C.m$} is the minimal set such that
  \begin{itemize}
  \item if $\kappa\ge 1$ and $v\in\dom(\tau_{C.m})$,
    then $v\in\paths{C.m}$, it has type $(\tau_{C.m})(v)$
    and $\length(v)=1$ (all local variables of $C.m$ are a path in $C.m$);
  \item if $\kappa\ge 1$ and $c\in\nat$ then $c\in\paths{C.m}$, it has type \<int>
    and $\length(c)=1$ (all integer constants are a path in $C.m$);
  \item if $\pi\in\paths{C.m}$ has type $C$, if $\length(\pi)<\kappa$
    and if $\mathtt{f}$ is a \<final> field of type $t$ defined in class $C$,
    then $\pi.f\in\paths{C.m}$, it has type $t$ and $\length(\pi.f)=\length(\pi)+1$
    (paths can be expanded by \<final> field dereference);
  \item if $\pi\in\paths{C.m}$ has type \<List<$t$\texttt{>}> and $\length(\pi)<\kappa$
    then $\pi\<[any]>\in\paths{C.m}$, it has type $t$ and $\length(\pi\<[any]>)=\length(\pi)+1$
    (paths can be expanded by list access);
  \item if $\pi\in\paths{C.m}$ has type \<List<$t$\texttt{>}> and $\length(\pi)<\kappa$
    then $\pi\<.size()>\in\paths{C.m}$, it has type $\<int>$ and $\length(\pi\<.size()>)=\length(\pi)+1$
    (paths can be expanded by list size computation).
  \end{itemize}
  A path is \emph{relevant} if its type can be assigned \<List<$t$\texttt{>}> for some $t$.
  The finite set of all relevant paths in $C.m$ is $\relevantpaths{C.m}$.
  When clear from the context, we write $\paths{}$ for $\paths{C.m}$ and
  $\relevantpaths{}$ for $\relevantpaths{C.m}$.
\end{definition}
%
For instance, at line~\ref{ponzi1:require} in Fig.~\ref{fig:ponzi1},
$42$ is a path of type \<int>, \<this> is a path of type \<Ponzi1>,
\<amount> is a path of type \<int>,
\<this.investors> is a path of type \<List<Contract\texttt{>}>,
\<this.investors[any]> is a path of type \<Contract>
and, finally, \<this.investors.size()> is a path of type \<int>. Among them, only
\<this.investors> is relevant.

A path \emph{survives} the execution of a command\ie it keeps its value unchanged,
if it does not access any list that might be modified by that execution. Note how this
definition is simplified by the hypothesis that local variables are not reassigned
(Def.~\ref{def:program}) and that paths only dereference
\<final> fields (Def.~\ref{def:path}). This allows one to focus on side-effects on lists only.
%
\begin{definition}\label{def:survival}
  Let $\pi\in\paths{C.m}$. Given a command $\com$ occurring in $C.m$, then
  \emph{$\pi$ survives $\com$}, written $\pi\survives{\com}$, if, for all prefix $\pi'$ of $\pi$
  (\ieinitial  $\pi=\pi'\cdots$) of type \<List<$t'$\texttt{>}>
  and for all $t''\in\sideeffect{\com}$, the types
  $t'$ and $t''$ have no common subtype.
  This survival test is extended to a filter on sets $S$ of paths, as $S\survives\com=\{\pi\in S\mid\pi\survives{\com}\}$,
  and further extended pointwise to functions $\sigma$ into sets of paths as
  $(\sigma\survives\com)(y)=\sigma(y)\survives\com$ for every $y\in\dom(\sigma)$. 
\end{definition}
%
For instance, the path $\pi=\<this.investors>$ at
line~\ref{ponzi1:expand} in Fig.~\ref{fig:ponzi1} does not survive
the command \<investors.push(caller)> at that line, since $\pi'=\<this.investors>$
has type \<List<Contract\texttt{>}>,
$\{\<Contract>\}=\sideeffect{\<investors.push(caller)>}$
and \<Contract> has a common subtype with itself, obviously.
Instead, the path $\pi=\<last.sales>$
survives the command \<rounds.push(last)>
at line~\ref{lottery2:push_round}, since
the only prefix $\pi'$ of $\pi$ of type list is $\pi'=\pi$, of type
\<List<TicketBatchSale\texttt{>}>,
$\{\<Round>\}=\sideeffect{\<rounds.push(last)>}$ and
\<TicketBatchSale> and \<Round> have no common subtype.
All integer constants survive any expression and command.

The definite aliasing analysis propagates abstract states that bind each
variable $v$ in scope to a set of paths that are definitely alias to $v$
(\ieinitial they have the same value as $v$).
%
\begin{definition}[Definite Aliasing Analysis States]\label{def:definite_aliasing_domain}
  Given a method $C.m$, the set of the
  \emph{definite aliasing states in $C.m$} is
  \[
  \dastates{C.m}=\{\sigma\mid\sigma:\dom(\tau_{C.m})\mapsto\paths{C.m}\}
  \]
  where $C.m$ is dropped when clear from the context.
\end{definition}
%
Since the definite aliasing analysis approximates states, it is more naturally
presented as an operational abstract semantics.
%
\begin{definition}[Definite Aliasing Analysis]\label{def:definite_aliasing}
  Given a method $C.m$, the \emph{definite aliasing semantics} of expressions
  in $C.m$ is given as a derivation function
  \[
  \left(\Exp\times\dastates{C.m}\right)\da{C.m}\left(\dastates{C.m}\times\paths{C.m}\right)
  \]
  from the approximation of the pre-states to the approximation of the post-states and
  the approximation of the paths aliased to the expression, where
  $C.m$ is dropped when clear from the context.
  Fig.~\ref{fig:daa_expressions} defines $\da{}$ for expressions.
%
  \begin{figure}[t]
    \begin{align*}
      \pair{c}{\sigma}&\da{}\pair{\sigma}{\{c\}\cap\paths{}}\\
      \pair{v}{\sigma}&\da{}\pair{\sigma}{(\sigma(v)\cup\{v\})\cap\paths{}}\\
      \pair{\field{v}{f}}{\sigma}&\da{}\pair{\sigma}{\{\pi.f\mid\pi\in\sigma(v)\}\cap\paths{}}\\
      \pair{\aop{v_1}{v_2}}{\sigma}&\da{}\pair{\sigma}{\{\aop{c_1}{c_2}\mid c_1\in\sigma(v_1),\ c_2\in\sigma(v_2)\}\cap\paths{}}\\
      \pair{\readmap{v_1}{v_2}{v_3}}{\sigma}&\da{}\pair{\sigma}{\varnothing}\\
      \pair{\new{C}{v^*}}{\sigma}&\da{}\pair{\sigma\survives{\new{C}{v^*}}}{\varnothing}\\
      \pair{\size{v}}{\sigma}&\da{}\pair{\sigma}{\{\pi\<.size()>\mid\pi\in\sigma(v)\}\cap\paths{}}\\
      \pair{\last{v}}{\sigma}&\da{}\pair{\sigma}{\{\pi\<[any]>\mid\pi\in\sigma(v)\}\cap\paths{}}\\
      \pair{\random{v}}{\sigma}&\da{}\pair{\sigma}{\varnothing}
    \end{align*}
    \caption{The definite aliasing derivation rules for expressions.}
    \label{fig:daa_expressions}
  \end{figure}

  The \emph{definite aliasing semantics} of a command $\com$ occurring in $C.m$ is
  a derivation function
  \[
  \left(\Com\times\dastates{C.m}\right)\da{C.m}\dastates{C.m}
  \]
  from the approximation of the pre-states to the approximation of the post-states,
  where $C.m$ is dropped when clear from the context.
  Fig.~\ref{fig:daa_commands} defines $\da{}$ for commands.
%
  \begin{figure}[t]
    \begin{gather*}
      \frac{\pair{\exp}{\sigma}\da{}\pair{\sigma'}{\val}}{\pair{\assign{v}{\exp}}{\sigma}\da{}\sigma'[v\mapsto\val]}\\[4pt]
      \pair{\assign{\field{v}{f}}{v}}{\sigma}\da{}\sigma\\[4pt]
      \pair{\writemap{v}{v}{v}}{\sigma}\da{}\sigma\\[4pt]
      \pair{\require{v}{v}}{\sigma}\da{}\sigma\\[4pt]
      \frac{\begin{array}{c}
          \sigma'=\sigma\survives{\push{v_1}{v_2}}\\
          \sigma''=\sigma'[v\mapsto\sigma'(v)\cup(\{\pi\<[any]>\mid\pi\in\sigma'(v_1)\}\cap\paths{})\mid\forall v\in\sigma(v_2)]
      \end{array}}{\pair{\push{v_1}{v_2}}{\sigma}\da{}\sigma''}\\[4pt]
      \pair{\receive{v}{v}}{\sigma}\da{}\sigma\\[4pt]
      \pair{\call{v}{m}{v^*}}{\sigma}\da{}\sigma\survives{\mathit{\call{v}{m}{v^*}}}\\[4pt]
      \frac{\pair{\com_1}{\sigma}\da{}\sigma'}{\pair{\ife{v}{v}{\com_1}{\com_2}}{\sigma}\da{}\sigma'}\\[4pt]
      \frac{\pair{\com_2}{\sigma}\da{}\sigma'}{\pair{\ife{v}{v}{\com_1}{\com_2}}{\sigma}\da{}\sigma'}\\[4pt]
      \frac{\pair{\com^*}{\sigma[v_1\mapsto\{\pi\<[any]>\mid\pi\in\sigma(v_2)\}\cap\paths{}]}\da{}\sigma'}{\pair{\for{v_1}{v_2}{\com}}{\sigma}\da{}\sigma'}\\[4pt]
      \pair{\com^*}{\sigma}\da{}\sigma\\[4pt]
      \frac{\pair{\com}{\sigma}\da{}\sigma'\quad\pair{\com^*}{\sigma'}\da{}\sigma''}{\pair{\com^*}{\sigma}\da{}\sigma''}\\[4pt]
      \frac{\pair{\com_1}{\sigma}\da{}\sigma'\quad\pair{\com_2}{\sigma'}\da{}\sigma''}{\pair{\sequence{\com_1}{\com_2}}{\sigma}=\sigma''}
    \end{gather*}
    \caption{The definite aliasing analysis for commands.}
    \label{fig:daa_commands}
  \end{figure}
%
\end{definition}
%
In Fig.~\ref{fig:daa_expressions},
the intersection with $\paths{}$ ensures that no illegal path
is generated (\ieinitial dereferencing non-\<final> fields or longer than $\kappa$).

The derivations in Fig.~\ref{fig:daa_commands} allow one to compute a definite
aliasing approximation at the beginning of every command in the program, through an
intraprocedural fixpoint that saturates all execution paths. Namely, one starts at
the beginning of a method with the less precise approximation, that maps each variable
to the empty set of paths. This approximation is propagated according to the rules
in Fig.~\ref{fig:daa_commands}, recording the abstract state whevener a new program point is
reached. If the propagation passes over an already visited program point, the
recorded approximation is the (per variable) intersection of the previously recorded
approximation and of the new approximation. This process must terminate since $\paths{}$ is finite
and intersection can only reduce the number of paths approximating each given variable.
In the following section, we assume that such definite aliasing approximation has been computed
at each command $\com$ in the program, and we write it as $\dastates{}@\com$.

